% 2012 年安徽卷文科第 12 题：直棱柱的三视图（单位：cm，图中只标数值）。
% 来源 PDF：普通高考/2012/2012安徽文.pdf，第 1 页（印面 715）中栏。
% 正视图 5×4，内部竖线在 x=2；侧视图 4×4；俯视图为直角梯形，上底 5、下底 2、高 4。
\begin{tikzpicture}[scale=0.70, semithick,
  dim/.style={<->, >=Stealth, thin},
  lab/.style={font=\small}]
  % ---------- 正视图 ----------
  \begin{scope}
    \draw (0,0) rectangle (5,4);
    \draw (2,0) -- (2,4);
    \draw[dim] (-0.45,0) -- (-0.45,4);
    \node[left] at (-0.45,2) {4};
    \node[below] at (2.5,-0.45) {正（主）视图};
  \end{scope}
  % ---------- 侧视图 ----------
  \begin{scope}[shift={(7.6,0)}]
    \draw (0,0) rectangle (4,4);
    \node[below] at (2,-0.45) {侧（左）视图};
  \end{scope}
  % ---------- 俯视图 ----------
  \begin{scope}[shift={(1.2,-6.6)}]
    \draw (0,0) -- (2,0) -- (5,4) -- (0,4) -- cycle;
    \draw[dim] (5.45,0) -- (5.45,4);
    \node[right] at (5.45,2) {4};
    \draw[dim] (0,4.4) -- (5,4.4);
    \node[above] at (2.5,4.4) {5};
    \draw[dim] (0,-0.4) -- (2,-0.4);
    \node[below] at (1,-0.4) {2};
    \node[below] at (2.5,-1.05) {俯视图};
  \end{scope}
\end{tikzpicture}
